\subsection{Auxiliary root pointers are erasable metadata (MODELBRIDGE-0020)}
\label{subsec:modelbridge0020}
\paragraph{Status.}
\textbf{PROVED/EXACT} for self-delimiting pointed cylinder encodings; P versus NP remains \textbf{OPEN}.

For a language $L$, set
\[
L^{\uparrow}=\{\langle x,u\rangle:x\in L,\ u\in\{0,1\}^*\}.
\]
Every uniform exact decider $A$ for $L$ induces a uniform exact decider
$\operatorname{Erase}[A]$ for $L^{\uparrow}$: parse and project $x$, discard
$u$, run $A(x)$, and return its one-bit answer.  With total pointed length
$M$ and $|x|=N$, the conservative sequential bounds are
\[
T_{\operatorname{Erase}[A]}(M)\le O(M)+T_A(N),\qquad
S_{\operatorname{Erase}[A]}(M)\le S_A(N)+O(\log M).
\]
Its post-projection computation is independent of $u$.  Consequently, a
compiler-designated root pointer that is only auxiliary provenance is not a
resource that every SAT execution must expose.  Combining this erasure lemma
with replacement-covariance factorization yields the exact interface audit:
auxiliary pointers are erasable; fully replacement-invariant pointers are
source data; implementation-dependent pointers do not support a universal
all-decider claim.  The surviving \textbf{OPEN} target is a source-defined,
predecision, strict hereditary compiler with a model-matched unconditional
endpoint after all auxiliary provenance is deleted.  The finite audit made
49,056 checks across exact-$(2,2)$, Horn, 2SAT, affine, bounded-treewidth,
product, padding, and UNSAT controls, with zero failures.
